% --- CAPÍTOL 2: MARC TEÒRIC ---
\chapter{Marc Teòric}

\section{Comunicació}
La comunicació és la base de la interacció humana. Segons la recerca realitzada, es divideix en diferents tipus:

\subsection{Verbal} 
Ús de la paraula (oral o escrita) per transmetre un missatge. És el mètode més comú.

\subsection{No verbal}
Són els missatges que enviem sense paraules: gestos, postura, contacte visual i expressions facials. En la comunicació diària, aquest tipus sol transmetre més informació emocional que les mateixes paraules.

\subsection{Comunicació Augmentativa i Alternativa (CAA)}
Sistemes, eines i estratègies que complementen (augmentativa) o substitueixen (alternativa) la parla quan aquesta no és funcional. Es divideix en dues branques:

\subsubsection{CAA Sense Ajuda (No assistida)}
No requereix cap objecte extern. La persona fa servir el seu propi cos.
{Exemples:} Llengua de signes, gestos codificats o l'ús de la mirada.

\subsubsection{CAA Amb Ajuda (Assistida):} Requereix l'ús d'una eina o suport extern.
{Baixa tecnologia:} Quaderns de comunicació, pictogrames (com el sistema SPC o ARASAAC) o llibres de fotos.
{Alta tecnologia:} Són aparells digitals que ajuden a parlar. Per exemple, aplicacions per a tauletes (com Proloquo2Go), equips que transformen el que escrius en veu o sistemes que permeten controlar l'ordinador només amb la mirada.



\section{Sistemes de suport}
S'han identificat diferents tipus de sistemes:
\begin{itemize}
    \item \textbf{Suport gràfic:} Ús de pictogrames. Destaca el portal \textbf{ARASAAC}, que és el portal de pictogrames més utilitzat.
    \item \textbf{Tecnologia:} Sistemes amb tecnologia i sistemes sense tecnologia.
\end{itemize}

\section{Integració sensorial}
Definició de com el cervell processa els estímuls. S'analitza la relació entre l'estímul sensorial i la resposta comunicativa, així com els beneficis de l'estimulació multisensorial.

\section{Dificultats comunicatives}
Aquest apartat analitza les causes (tipus), les barreres socials i l'impacte emocional que pateixen els usuaris. La taula sensorial té una utilitat real per a:
\begin{itemize}
    \item TEA (Trastorn de l'espectre autista).
    \item Paràlisi cerebral.
    \item Infants amb retard del llenguatge.
    \item Discapacitat intel·lectual.
\end{itemize}
